% Caso de Uso: Agregar Documento
% Sistema: Hotel Perros

%-------------------------------------- COMIENZA descripción del caso de uso.

\begin{UseCase}{CU12}{Agregar Documento}{
    Permite al Propietario subir y asociar un archivo digital (imagen, PDF o documento de texto) al perfil de su mascota, especificando opcionalmente el tipo de documento.
}
    \UCitem{Versión}{\color{Gray}1.0}
    \UCitem{Autor}{\color{Gray}Ugalde Tellez Aaron}
    \UCitem{Supervisa}{\color{Gray}Castillo Aldaba Jared}
    \UCitem{Actor}{\hyperlink{Propietario}{Propietario}}
    \UCitem{Propósito}{Almacenar documentación relevante de la mascota (cartilla de vacunación, pedigree, resultados médicos) para su consulta digital centralizada.}
    \UCitem{Entradas}{
        \begin{itemize}
            \item Selección del Tipo de Documento (Opcional).
            \item Archivo digital (Obligatorio: PDF, DOC, DOCX o Imagen hasta 10MB).
        \end{itemize}
    }
    \UCitem{Origen}{Mouse / Pantalla táctil y Sistema de Archivos Local}
    \UCitem{Salidas}{Mensaje de confirmación (implícito al redirigir), visualización del documento en el perfil de la mascota.}
    \UCitem{Destino}{Pantalla}
    \UCitem{Precondiciones}{
        \begin{itemize}
            \item El propietario ha iniciado sesión.
            \item El propietario ha seleccionado una mascota y se encuentra en su perfil (CU02 Consultar Perfil de Mascota).
        \end{itemize}
    }
    \UCitem{Postcondiciones}{El archivo es almacenado en el servidor y se crea un registro en la base de datos asociado a la mascota.}
    \UCitem{Errores}{
        \begin{itemize}
            \item \textbf{E1}: El formato del archivo no es válido.
            \item \textbf{E2}: El tamaño del archivo excede el límite permitido.
            \item \textbf{E3}: Error de conexión al subir el archivo.
        \end{itemize}
    }
    \UCitem{Tipo}{Secundario}
\end{UseCase}

%--------------------------------------
% Trayectoria Principal
%--------------------------------------
\begin{UCtrayectoria}
    \UCpaso[\UCactor] Oprime el botón \IUbutton{Agregar Documento} desde la vista de detalles de la mascota.
    \UCpaso Muestra la pantalla "Subir Documento" con:
        \begin{itemize}
            \item Lista desplegable de Tipos de Documento.
            \item Área de carga de archivos (Drag \& Drop o selección).
            \item Botones \IUbutton{Cancelar} y \IUbutton{Subir Documento} (deshabilitado inicialmente).
        \end{itemize}
    \UCpaso[\UCactor] Selecciona una opción en la lista "Tipo de Documento" (Ej. Certificado veterinario, Cartilla, etc.) o lo deja "Sin especificar".
    \UCpaso[\UCactor] Carga un archivo haciendo clic en "Seleccionar archivo" o arrastrándolo al área designada.
    \UCpaso Muestra el nombre del archivo seleccionado y habilita el botón \IUbutton{Subir Documento}.
    \UCpaso[\UCactor] Oprime el botón \IUbutton{Subir Documento}.
    \UCpaso Valida que se haya seleccionado un archivo válido.
    \UCpaso Sube el archivo al servidor y guarda la referencia en la base de datos.
    \UCpaso Redirige a la pantalla de detalles de la mascota mostrando el nuevo documento en la lista.
\end{UCtrayectoria}

%--------------------------------------
% Trayectorias Alternativas
%--------------------------------------

% Trayectoria A: Cancelar operación
\begin{UCtrayectoriaA}{A}{El propietario cancela la operación}
    \UCpaso[\UCactor] Oprime el botón \IUbutton{Cancelar} o el enlace "Volver a [Nombre Mascota]".
    \UCpaso Regresa a la pantalla de detalles de la mascota sin guardar ningún documento.
\end{UCtrayectoriaA}

% Trayectoria B: Intento de subir sin archivo
\begin{UCtrayectoriaA}{B}{El usuario intenta subir sin seleccionar archivo}
    \UCpaso Detecta que la variable del archivo está vacía.
    \UCpaso Muestra el mensaje de error {\bf MSG-Error} ``Por favor selecciona un archivo''.
    \UCpaso Mantiene al usuario en el formulario para que pueda corregir.
\end{UCtrayectoriaA}

% Trayectoria C: Error del servidor
\begin{UCtrayectoriaA}{C}{Fallo en la carga al servidor}
    \UCpaso Intenta subir el documento pero recibe una respuesta de error del API.
    \UCpaso Muestra el mensaje de error {\bf MSG-Error} (Ej. ``Error al subir el documento'' o mensaje específico del servidor).
    \UCpaso Continúa en el paso 2 de la trayectoria principal permitiendo reintentar.
\end{UCtrayectoriaA}